\section{Discussion}\label{sec:discussion}

A builder of web agents can take four things from this without adopting our conclusion.

\looseness=-1 \textbf{Measure your own cell.} Which representation wins is a property of the deployment, not of the modality: it reverses between our two task sets and it is not stable under rerun noise in the low-success cells. A number read off someone else's benchmark does not transfer to your site, and the eight cells here disagree with each other more than any of them disagrees with a plausible prior.

\looseness=-1 \textbf{Price a rerun before pricing a representation.} The cheapest way to raise a union ceiling is to run an arm you already have a second time. Until that is measured, a gain from adding a representation cannot be attributed to the representation. Our own headline ceiling survived as a measurement but not as an attribution; the cost ceiling survived as both, because it adds no arm.

\textbf{Treat a difference smaller than the rerun band as unresolved, not as small.} At the flip rate we observe, a two-point difference between two modes is what a single repetition delivers on its own. Reporting it as an effect is reporting the resolution of the instrument.

\looseness=-1 \textbf{State the denominator, and do not pool a bill with an energy estimate.} Per-attempt and per-success cost disagree about which mode is cheapest in most of our cells; an API bill and an electricity estimate are not one quantity; and a latency figure is mostly browser and container. Each silently reorders the modes, so an efficiency claim without its estimand is ambiguous rather than weak: the trouble is not that the effect is small, it is that one cannot tell what was measured.

\looseness=-1 None of this settles whether representation routing can be made to work. It bounds the question: what a perfect choice would buy, what the available supervision delivers, and how far apart the two sit on the workloads the field measures.
